%================================================================%
%   DRAFT: Section 3 — Copula Generated Random Graphs            %
%                       for Real-Valued Networks                 %
%                                                                %
%   Insert between the binary CGRG section and the brain        %
%   application section in CGRG_brain.tex                       %
%================================================================%

\section{Copula Generated Random Graphs for Real-Valued Networks}
\label{sec:realvalued}

Many networks arising in the natural and social sciences carry real-valued
edge weights.  Regulatory strength between brain regions, trade volumes
between countries, and communication frequencies between colleagues are
all examples.  This section develops a Copula Generated Random Graph
model for directed networks with real-valued edges, generalising the
binary case of Section~\ref{sec:binary}.

\subsection{Model Specification}

Let $G$ be a directed network on $N$ nodes with edge weight matrix
$A \in \mathbb{R}^{N \times N}$.  Following the block-separable construction
of Section~\ref{sec:cgrg}, we partition the off-diagonal edges
into $M$ groups such that edges belonging to different groups are
conditionally independent.  In what follows we take the natural partition
into dyads, so $M = N(N-1)/2$ and each group contains the two directed
edges $A_{ij}$ and $A_{ji}$ connecting the same pair of nodes.  Under this
partition, the log-likelihood of the network decomposes as
%
\begin{align}
\label{eq:real_loglik}
\log p(A \mid \theta, \rho)
  = \sum_{i=1}^{N} \log p(A_{ii} \mid \theta)
  + \sum_{i=1}^{N-1} \sum_{j > i}
    \log p(A_{ij}, A_{ji} \mid \theta, \rho),
\end{align}
%
where $\theta$ collects the marginal parameters and $\rho$ is the copula
dependence parameter.

\subsection{Marginal Distributions}

The marginal distribution of each off-diagonal edge $A_{ij}$, $i \neq j$,
is specified as
%
\begin{align}
\label{eq:marginal}
A_{ij} \mid \delta_i, \gamma_j, \lambda
  \;\sim\; \mathcal{N}\!\left(\delta_i + \gamma_j,\; \lambda^{-1}\right),
\end{align}
%
where $\delta_i$ is a \emph{sender effect} for node $i$, measuring the
tendency of node $i$ to exert influence on other nodes, $\gamma_j$ is a
\emph{receiver effect} for node $j$, measuring the tendency of node $j$ to
be influenced by other nodes, and $\lambda > 0$ is a precision parameter
shared across all edges in the network.  The mean structure $\delta_i +
\gamma_j$ is a standard sender-receiver additive model \citep{hoff2005}.

\paragraph{Identifiability.}
The sender and receiver effects are not separately identified from the
overall level of $A$ without a constraint.  We fix $\delta_1 = 0$,
so that all sender effects are interpreted relative to node~$1$.  With
this normalisation, $\gamma_j$ measures the unconditional mean weight of
edges directed into node $j$, and $\delta_i$ for $i \geq 2$ measures
the excess sending tendency of node $i$ relative to the reference node.

\subsection{Frank Copula Density for a Dyad}

By Sklar's theorem \citep{sklar1959}, the joint distribution of the two
edges $A_{ij}$ and $A_{ji}$ can be expressed via a copula $C_\rho$
binding their marginal CDFs $F_{ij}$ and $F_{ji}$:
%
\begin{align}
p(A_{ij}, A_{ji} \mid \theta, \rho)
  = c_\rho\!\left(F_{ij}(A_{ij}), F_{ji}(A_{ji})\right)
    \cdot f_{ij}(A_{ij}) \cdot f_{ji}(A_{ji}),
\end{align}
%
where $c_\rho(u,v) = \partial^2 C_\rho(u,v) / \partial u\, \partial v$
is the copula density.  We choose the bivariate \emph{Frank copula}
%
\begin{align}
\label{eq:frank_cdf}
C_\rho(u,v)
  = -\frac{1}{\rho}
    \log\!\left[1 + \frac{(e^{-\rho u}-1)(e^{-\rho v}-1)}{e^{-\rho}-1}\right],
\qquad \rho \in \mathbb{R} \setminus \{0\},
\end{align}
%
whose density is
%
\begin{align}
\label{eq:frank_density}
c_\rho(u,v)
  = \frac{\rho(1 - e^{-\rho})\, e^{-\rho(u+v)}}
         {\left[1 - e^{-\rho} - (1-e^{-\rho u})(1-e^{-\rho v})\right]^2}.
\end{align}
%
With Gaussian marginals $F_{ij}(a) = \Phi\!\left((a - \delta_i - \gamma_j)
\sqrt{\lambda}\right)$, the joint log-likelihood contribution of dyad $ij$
is
%
\begin{align}
\label{eq:dyad_loglik}
\log p(A_{ij}, A_{ji} \mid \theta, \rho)
  &= \log c_\rho\!\left(\Phi(z_{ij}),\, \Phi(z_{ji})\right)
  + \log \phi(z_{ij}) + \tfrac{1}{2}\log\lambda  \nonumber\\
  &\quad + \log\phi(z_{ji}) + \tfrac{1}{2}\log\lambda,
\end{align}
%
where $z_{ij} = (A_{ij} - \delta_i - \gamma_j)\sqrt{\lambda}$,
$\phi$ denotes the standard normal density, and the copula density is
evaluated at the probability-integral-transformed edge values.

\paragraph{The role of $\rho$.}
The Frank copula dependence parameter $\rho \in \mathbb{R}\setminus\{0\}$
controls the degree and direction of \emph{reciprocity} in the network.
When $\rho > 0$, edge weights $A_{ij}$ and $A_{ji}$ are positively
associated: a node that sends a strong signal to its partner tends to
receive a strong signal in return.  When $\rho < 0$, the association is
negative: strong outgoing edges tend to coincide with weak incoming edges,
implying a hierarchical or asymmetric regulatory structure.  As
$\rho \to 0$, the two directed edges become independent.  This bidirectional
interpretability distinguishes the Frank copula from, e.g., the
Gaussian copula, while retaining a closed-form density and straightforward
computation.

\subsection{Full Likelihood}

Combining Equations~\eqref{eq:real_loglik}
and~\eqref{eq:dyad_loglik}, the log-likelihood for a single
network $A$ is
%
\begin{align}
\label{eq:full_loglik}
\log p(A \mid \theta, \rho)
  &= \sum_{i=1}^{N} \log p(A_{ii} \mid \theta) \nonumber \\
  &\quad+
    \sum_{i=1}^{N-1}\sum_{j>i}
    \Bigl[
      \log c_\rho\!\left(\Phi(z_{ij}),\Phi(z_{ji})\right)
      + \log\phi(z_{ij}) + \log\phi(z_{ji}) + \log\lambda
    \Bigr].
\end{align}
%
When $S$ independent network replicates $A^{(1)},\ldots,A^{(S)}$ are
observed, the log-likelihood is the sum of~\eqref{eq:full_loglik} over
replicates.

\subsection{Bayesian Inference via Markov Chain Monte Carlo}

We adopt a Bayesian framework.  The hierarchical prior specification is:
%
\begin{align}
\delta_0, \gamma_0 &\;\sim\; \mathcal{N}(0, 4),
  \label{eq:prior_hyper} \\
\rho_0 &\;\sim\; \mathcal{N}(0, 4), \\
\delta_i \mid \delta_0, \sigma_\delta
  &\;\sim\; \mathcal{N}(\delta_0, \sigma_\delta^2),
  \quad i = 2,\ldots,N, \\
\gamma_j \mid \gamma_0, \sigma_\gamma
  &\;\sim\; \mathcal{N}(\gamma_0, \sigma_\gamma^2),
  \quad j = 1,\ldots,N, \\
\rho \mid \rho_0, \sigma_\rho
  &\;\sim\; \mathcal{N}(\rho_0, \sigma_\rho^2), \\
\lambda &\;\sim\; \mathrm{Gamma}(5,\, 4), \\
\sigma_\rho, \sigma_\delta, \sigma_\gamma
  &\;\sim\; \mathrm{Gamma}(0.1,\, 10).
  \label{eq:prior_sigma}
\end{align}
%
The hyper-mean parameters $(\delta_0, \gamma_0, \rho_0)$ share
information across nodes and groups respectively, while the
hyper-variance parameters $(\sigma_\delta, \sigma_\gamma, \sigma_\rho)$
control the degree of shrinkage.  The Gamma prior on $\lambda$
concentrates mass on precisions consistent with unit-order edge weights,
while the weakly-informative Gamma priors on the variance hyperparameters
allow the data to determine the extent of pooling.

Posterior samples are drawn using a componentwise
Metropolis--Hastings algorithm in which each parameter is updated in
turn using a symmetric uniform proposal.  The posterior is
%
\begin{align}
p(\theta, \rho \mid A^{(1)},\ldots,A^{(S)})
  \;\propto\;
  \left[\prod_{s=1}^{S} p(A^{(s)} \mid \theta, \rho)\right]
  p(\theta, \rho),
\end{align}
%
where $p(\theta,\rho)$ is the joint prior implied by
Equations~\eqref{eq:prior_hyper}--\eqref{eq:prior_sigma}.
The MCMC sampler cycles through $\rho$, $\{\delta_i\}$,
$\{\gamma_j\}$, $\lambda$, and the hyperparameters in sequence,
evaluating the log posterior via~\eqref{eq:full_loglik} at each
step.  We discard an initial burn-in and thin the chain to obtain
approximately independent draws for posterior summary.

\subsection{Relationship to the Multivariate Normal Model}

A natural comparison model treats the collection of $N(N-1)$ off-diagonal
edge weights for each network as a multivariate normal vector with an
unstructured covariance matrix.  While flexible, this model has
$O(N^4)$ covariance parameters, which is difficult to estimate with
typical sample sizes, and offers little structural interpretability.
The CGRG model achieves parsimony by restricting dependence to
within-dyad pairs — reducing the dependence structure to a single
scalar $\rho$ — while preserving the interpretable sender/receiver
decomposition of the mean.  Specifically, under the CGRG model
the marginal covariance between $A_{ij}$ and $A_{ji}$ is determined
entirely by $\rho$ through the Frank copula, whereas all pairs of edges
from different dyads are modeled as independent given the node effects.
The KL divergence between the CGRG density and the bivariate normal
density for a single dyad quantifies the misspecification incurred
when using the multivariate normal in place of the copula model;
this comparison is developed in Section~\ref{sec:compare}.

% ---------------------------------------------------------------- %
%  NOTES FOR REVISION:                                             %
%   - Add citation to Hoff (2005) for sender/receiver model       %
%   - Add citation to Sklar (1959) / Nelsen (2006) for copulas    %
%   - Add citation to Joe (1997) for Frank copula properties      %
%   - Consider adding a remark about vine copula extension        %
%     (see Czado 2025, CSDA) for more than pairwise dependence    %
%   - The Feng & Leng (Biometrika 2025) paper establishes         %
%     consistency/asymptotics for a related directed network model %
%     with sender/receiver + reciprocity -- worth citing          %
% ---------------------------------------------------------------- %
